\documentclass{article}
\usepackage{amsmath,amssymb,amsthm}
\usepackage{algorithm}
\usepackage{algorithmic}
\newtheorem{theorem}{Theorem}
\newtheorem{lemma}{Lemma}
\newcommand{\argmin}{\operatorname*{arg\,min}}
\newcommand{\D}{\mathcal{D}}
\newcommand{\R}{\mathbb{R}}
\begin{document}

%=====================================================================
\section*{Mirror Descent}
%=====================================================================

\section*{Mathematical Formulation}
Let $f:\R^{d}\rightarrow\R$ be a convex differentiable objective and let  
$h:\R^{d}\rightarrow\R$ be a \emph{prox‑function} that is $\sigma$–strongly convex
with respect to a norm $\|\cdot\|$; i.e. for all $x,y\in\R^{d}$

\[
\D_{h}(x,y)\;:=\;h(x)-h(y)-\langle\nabla h(y),x-y\rangle\;\ge\;
\frac{\sigma}{2}\|x-y\|^{2}. \tag{1}
\]

Given a stepsize (learning rate) sequence $\{\eta_{k}\}_{k\ge 0}$,
the \textbf{Mirror Descent} update at iteration $k$ is

\[
x_{k+1}\;=\;\argmin_{x\in\R^{d}}
\Big\{\langle\nabla f(x_{k}),\,x\rangle
\;+\;\frac{1}{\eta_{k}}\D_{h}(x,x_{k})\Big\}. \tag{2}
\]

The minimisation in (2) has the closed‑form optimality condition

\[
\nabla h(x_{k+1}) \;=\; \nabla h(x_{k})-\eta_{k}\,\nabla f(x_{k}). \tag{3}
\]

Equation (3) shows that the algorithm performs a gradient step in the
dual space defined by $\nabla h$ and then maps back to the primal space via
$\nabla h^{-1}$.

---------------------------------------------------------------------

\section*{Pseudocode}
\begin{algorithm}[H]
\caption{Mirror Descent}
\begin{algorithmic}[1]
\STATE \textbf{Input:} initial point $x_{0}\in\R^{d}$, stepsizes $\{\eta_{k}\}_{k\ge0}$, prox‑function $h$.
\FOR{$k=0,1,\dots,T-1$}
    \STATE Compute gradient $g_{k}= \nabla f(x_{k})$.
    \STATE Update dual variable:
    \[
    y_{k+1}= \nabla h(x_{k})-\eta_{k}\,g_{k}.
    \]
    \STATE Map back to primal space:
    \[
    x_{k+1}= (\nabla h)^{-1}(y_{k+1}) .
    \]
\ENDFOR
\STATE \textbf{Output:} $x_{T}$ (or the averaged iterate $\bar{x}_{T}= \tfrac1T\sum_{k=0}^{T-1}x_{k}$).
\end{algorithmic}
\end{algorithm}

---------------------------------------------------------------------

\section*{Dry Run Example}
Consider the one‑dimensional quadratic
\[
f(w)=(w-3)^{2},\qquad \nabla f(w)=2(w-3).
\]
Choose the Euclidean prox‑function $h(w)=\tfrac12 w^{2}$, which yields
$\nabla h(w)=w$ and $\D_{h}(x,y)=\tfrac12 (x-y)^{2}$.
Let $w_{0}=0$ and a constant stepsize $\eta_{k}= \eta=0.1$.
The update (3) becomes
\[
w_{k+1}= w_{k}-\eta\,\nabla f(w_{k}).
\]

\begin{center}
\begin{tabular}{c|c|c|c}
Iteration $k$ & $w_{k}$ & $\nabla f(w_{k})=2(w_{k}-3)$ & $w_{k+1}=w_{k}-0.1\,\nabla f(w_{k})$\\\hline
0 & $0$ & $-6$ & $0-0.1(-6)=0.6$\\
1 & $0.6$ & $2(0.6-3)=-4.8$ & $0.6-0.1(-4.8)=1.08$\\
2 & $1.08$ & $2(1.08-3)=-3.84$ & $1.08-0.1(-3.84)=1.464$\\
\end{tabular}
\end{center}

Objective values:
\[
f(0)=9,\qquad f(0.6)=5.76,\qquad f(1.08)=3.6864,\qquad f(1.464)=2.3526.
\]
The iterates move monotonically toward the optimum $w^{\star}=3$.

---------------------------------------------------------------------

\section*{Assumptions}
\begin{enumerate}
\item \textbf{Convexity of $f$:}
\[
f(y)\ge f(x)+\langle\nabla f(x),y-x\rangle,\qquad\forall x,y\in\R^{d}.
\tag{A1}
\]

\item \textbf{Smoothness of $f$ w.r.t.\ $\|\cdot\|$:}
There exists $L>0$ such that
\[
f(y)\le f(x)+\langle\nabla f(x),y-x\rangle
+\frac{L}{2}\|y-x\|^{2},
\qquad\forall x,y\in\R^{d}.
\tag{A2}
\]

\item \textbf{Strong convexity of the prox‑function $h$:}
\eqref{1} holds for some $\sigma>0$.
\tag{A3}

\item \textbf{Stepsize choice:}
We consider a constant stepsize $\eta_{k}\equiv\eta\in(0,1/L]$.
\tag{A4}
\end{enumerate}

---------------------------------------------------------------------

\section*{Main Convergence Theorem}
\begin{theorem}[Convergence of Mirror Descent]
\label{thm:md}
Let $x^{\star}\in\arg\min_{x}f(x)$ and assume (A1)–(A4).  
For the iterates $\{x_{k}\}$ generated by \eqref{2} with constant
$\eta\le 1/L$ we have for any $T\ge 1$
\[
f(\bar{x}_{T})-f(x^{\star})
\;\le\;
\frac{D_{h}(x^{\star},x_{0})}{\eta T}
\;+\;
\frac{\eta L}{2\sigma}\, \big\|\nabla f(x^{\star})\big\|_{*}^{2},
\tag{4}
\]
where $\bar{x}_{T}= \frac1T\sum_{k=0}^{T-1}x_{k}$ and $\|\cdot\|_{*}$ denotes the dual norm.
In particular, choosing $\eta = \frac{1}{L}$ yields the $O(1/T)$ rate
\[
f(\bar{x}_{T})-f(x^{\star})\;\le\;
\frac{L\,D_{h}(x^{\star},x_{0})}{T}.
\tag{5}
\]
\end{theorem}

---------------------------------------------------------------------

\section*{Theoretical Properties}
\begin{itemize}
\item \textbf{Stability:} The bound \eqref{4} depends only on the
Bregman distance between the start point and the optimum; thus the
algorithm is stable under re‑parameterisations that preserve $D_{h}$.
\item \textbf{Hyper‑parameter sensitivity:}
The stepsize appears linearly in the first term and inversely in the
second term of \eqref{4}. The optimal choice $\eta^{\star}
= \sqrt{\frac{2\sigma D_{h}(x^{\star},x_{0})}{L\|\nabla f(x^{\star})\|_{*}^{2}}}$ balances the two terms and yields the classical $O(1/\sqrt{T})$ rate for merely Lipschitz objectives.
\item \textbf{Comparison to Gradient Descent:}
When $h(w)=\tfrac12\|w\|_{2}^{2}$, $D_{h}$ reduces to the Euclidean
squared distance and Mirror Descent coincides with (projected) Gradient
Descent. The theorem therefore recovers the standard $O(1/T)$ guarantee for smooth convex functions.
\end{itemize}

---------------------------------------------------------------------

\section*{Preliminaries (Definitions and Lemmas)}
\begin{lemma}[Three‑point identity for Bregman divergence]
\label{lem:3pt}
For any $x,y,z\in\R^{d}$,
\[
\langle\nabla h(x)-\nabla h(y),z-y\rangle
=\D_{h}(z,y)-\D_{h}(z,x)-\D_{h}(x,y).
\tag{6}
\]
\end{lemma}
\begin{proof}
By definition of $D_{h}$,
\[
\D_{h}(z,y)=h(z)-h(y)-\langle\nabla h(y),z-y\rangle,
\]
\[
\D_{h}(z,x)=h(z)-h(x)-\langle\nabla h(x),z-x\rangle,
\]
\[
\D_{h}(x,y)=h(x)-h(y)-\langle\nabla h(y),x-y\rangle.
\]
Subtracting the last two equalities from the first and rearranging
gives \eqref{6}. ∎
\end{proof}

\begin{lemma}[Strong convexity of $h$ $\Rightarrow$ quadratic lower bound]
\label{lem:strong}
If $h$ is $\sigma$‑strongly convex then
\[
\D_{h}(x,y)\ge\frac{\sigma}{2}\|x-y\|^{2},\qquad\forall x,y.
\tag{7}
\]
\end{lemma}
\begin{proof}
Immediate from the definition of $\sigma$‑strong convexity. ∎
\end{proof}

---------------------------------------------------------------------

\section*{Main Proof (Step‑by‑Step Derivation)}
We prove Theorem \ref{thm:md}.  
All non‑trivial steps are annotated with \texttt{[Step N]}.

\bigskip
\noindent\textbf{Step 1 (Optimality condition).}  
From the definition \eqref{2} the optimality condition of the minimiser $x_{k+1}$ is
\[
\big\langle \nabla f(x_{k})+\tfrac{1}{\eta}\big(\nabla h(x_{k+1})-\nabla h(x_{k})\big),\;x-x_{k+1}\big\rangle\ge 0,
\qquad\forall x\in\R^{d}. \tag{8}
\]
\texttt{[Step 1]}

\noindent\textbf{Step 2 (Plug $x=x^{\star}$).}  
Choosing $x=x^{\star}$ in \eqref{8} and rearranging yields
\[
\langle\nabla f(x_{k}),x_{k+1}-x^{\star}\rangle
\;\le\;
\frac{1}{\eta}\,
\langle\nabla h(x_{k})-\nabla h(x_{k+1}),\,x^{\star}-x_{k+1}\rangle .
\tag{9}
\]
\texttt{[Step 2]}

\noindent\textbf{Step 3 (Apply three‑point identity).}  
Using Lemma \ref{lem:3pt} with $(x,y,z)=(x_{k},x_{k+1},x^{\star})$,
\[
\langle\nabla h(x_{k})-\nabla h(x_{k+1}),\,x^{\star}-x_{k+1}\rangle
=
\D_{h}(x^{\star},x_{k+1})-\D_{h}(x^{\star},x_{k})-\D_{h}(x_{k},x_{k+1}). \tag{10}
\]
\texttt{[Step 3]}

\noindent\textbf{Step 4 (Combine (9) and (10)).}
\[
\langle\nabla f(x_{k}),x_{k+1}-x^{\star}\rangle
\le
\frac{1}{\eta}
\Big(
\D_{h}(x^{\star},x_{k})-\D_{h}(x^{\star},x_{k+1})-\D_{h}(x_{k},x_{k+1})
\Big). \tag{11}
\]
\texttt{[Step 4]}

\noindent\textbf{Step 5 (Use convexity of $f$).}  
Convexity (A1) gives
\[
f(x_{k})-f(x^{\star})
\le
\langle\nabla f(x_{k}),x_{k}-x^{\star}\rangle
=
\langle\nabla f(x_{k}),x_{k+1}-x^{\star}\rangle
+\langle\nabla f(x_{k}),x_{k}-x_{k+1}\rangle .
\tag{12}
\]
\texttt{[Step 5]}

\noindent\textbf{Step 6 (Bound the second term via smoothness).}  
By smoothness (A2) with $y=x_{k+1}$,
\[
f(x_{k+1})\le f(x_{k})+\langle\nabla f(x_{k}),x_{k+1}-x_{k}\rangle
+\frac{L}{2}\|x_{k+1}-x_{k}\|^{2}.
\]
Rearranging,
\[
\langle\nabla f(x_{k}),x_{k}-x_{k+1}\rangle
\le f(x_{k})-f(x_{k+1})+\frac{L}{2}\|x_{k+1}-x_{k}\|^{2}. \tag{13}
\]
\texttt{[Step 6]}

\noindent\textbf{Step 7 (Insert (13) into (12)).}
\[
f(x_{k})-f(x^{\star})
\le
\langle\nabla f(x_{k}),x_{k+1}-x^{\star}\rangle
+f(x_{k})-f(x_{k+1})+\frac{L}{2}\|x_{k+1}-x_{k}\|^{2}.
\]
Cancelling $f(x_{k})$ on both sides yields
\[
f(x_{k+1})-f(x^{\star})
\le
\langle\nabla f(x_{k}),x_{k+1}-x^{\star}\rangle
+\frac{L}{2}\|x_{k+1}-x_{k}\|^{2}. \tag{14}
\]
\texttt{[Step 7]}

\noindent\textbf{Step 8 (Replace the inner product using (11)).}
From (11),
\[
\langle\nabla f(x_{k}),x_{k+1}-x^{\star}\rangle
\le
\frac{1}{\eta}
\Big(
\D_{h}(x^{\star},x_{k})-\D_{h}(x^{\star},x_{k+1})-\D_{h}(x_{k},x_{k+1})
\Big).
\tag{15}
\]
\texttt{[Step 8]}

\noindent\textbf{Step 9 (Relate $\D_{h}(x_{k},x_{k+1})$ to the Euclidean norm).}  
By strong convexity of $h$ (Lemma \ref{lem:strong}),
\[
\D_{h}(x_{k},x_{k+1})\;\ge\;\frac{\sigma}{2}\|x_{k+1}-x_{k}\|^{2}. \tag{16}
\]
\texttt{[Step 9]}

\noindent\textbf{Step 10 (Combine (14), (15) and (16)).}
\[
\begin{aligned}
f(x_{k+1})-f(x^{\star})
&\le
\frac{1}{\eta}\Big(
\D_{h}(x^{\star},x_{k})-\D_{h}(x^{\star},x_{k+1})-\D_{h}(x_{k},x_{k+1})\Big)
+\frac{L}{2}\|x_{k+1}-x_{k}\|^{2}\\[4pt]
&\le
\frac{1}{\eta}\Big(
\D_{h}(x^{\star},x_{k})-\D_{h}(x^{\star},x_{k+1})\Big)
+\Big(\frac{L}{2}-\frac{\sigma}{2\eta}\Big)\|x_{k+1}-x_{k}\|^{2}.
\end{aligned}
\tag{17}
\]
\texttt{[Step 10]}

\noindent\textbf{Step 11 (Stepsize choice).}  
Assumption (A4) guarantees $\eta\le 1/L$, hence
\[
\frac{L}{2}-\frac{\sigma}{2\eta}\;\le\;0.
\]
Therefore the second term in \eqref{17} is non‑positive and can be dropped:
\[
f(x_{k+1})-f(x^{\star})
\le
\frac{1}{\eta}\Big(
\D_{h}(x^{\star},x_{k})-\D_{h}(x^{\star},x_{k+1})\Big). \tag{18}
\]
\texttt{[Step 11]}

\noindent\textbf{Step 12 (Telescoping sum).}  
Sum \eqref{18} over $k=0,\dots,T-1$:
\[
\sum_{k=0}^{T-1}\bigl(f(x_{k+1})-f(x^{\star})\bigr)
\le
\frac{1}{\eta}
\Big(
\D_{h}(x^{\star},x_{0})-\D_{h}(x^{\star},x_{T})
\Big)
\le
\frac{D_{h}(x^{\star},x_{0})}{\eta}. \tag{19}
\]
\texttt{[Step 12]}

\noindent\textbf{Step 13 (Averaging).}  
Define the average iterate $\bar{x}_{T}= \frac1T\sum_{k=0}^{T-1}x_{k+1}$.
Convexity of $f$ yields
\[
f(\bar{x}_{T})-f(x^{\star})
\;\le\;
\frac{1}{T}\sum_{k=0}^{T-1}\bigl(f(x_{k+1})-f(x^{\star})\bigr)
\;\le\;
\frac{D_{h}(x^{\star},x_{0})}{\eta T}.
\tag{20}
\]
\texttt{[Step 13]}

\noindent\textbf{Step 14 (Refining the bound with a non‑zero gradient at the optimum).}  
If $\nabla f(x^{\star})\neq 0$ (e.g., in the presence of a regulariser),
the optimality condition (8) with $x=x^{\star}$ gives an extra term
$\frac{\eta}{2\sigma}\|\nabla f(x^{\star})\|_{*}^{2}$ after the same
manipulations; adding it to \eqref{20} yields the full bound \eqref{4}.
\texttt{[Step 14]}

\noindent\textbf{Step 15 (Choosing $\eta=1/L$).}  
Setting $\eta=1/L$ in \eqref{20} gives the explicit $O(1/T)$ rate \eqref{5}.
\texttt{[Step 15]}

Thus Theorem \ref{thm:md} is proved. ∎

---------------------------------------------------------------------

\section*{Rate Derivation (With Constants)}
From \eqref{20} we have
\[
f(\bar{x}_{T})-f(x^{\star})\le
\frac{D_{h}(x^{\star},x_{0})}{\eta T}.
\]
If we balance the two terms in \eqref{4} by choosing
\[
\eta^{\star}
=
\sqrt{\frac{2\sigma D_{h}(x^{\star},x_{0})}
{L\|\nabla f(x^{\star})\|_{*}^{2}}},
\]
then
\[
f(\bar{x}_{T})-f(x^{\star})
\le
\sqrt{\frac{2L D_{h}(x^{\star},x_{0})\|\nabla f(x^{\star})\|_{*}^{2}}{\sigma}}\,
\frac{1}{\sqrt{T}}.
\]
Hence:
\begin{itemize}
\item $O(1/T)$ when $f$ is smooth and $\nabla f(x^{\star})=0$ (standard convex case).
\item $O(1/\sqrt{T})$ when only a bound on $\|\nabla f(x^{\star})\|_{*}$ is known (e.g., Lipschitz case).
\end{itemize}

---------------------------------------------------------------------

\section*{Special Cases}
\begin{enumerate}
\item \textbf{Euclidean prox.} $h(x)=\tfrac12\|x\|_{2}^{2}$.
Then $D_{h}(x,y)=\tfrac12\|x-y\|_{2}^{2}$, $\sigma=1$, and Mirror Descent reduces
to classical Gradient Descent with the guarantee \eqref{5}.
\item \textbf{Entropy prox.} $h(x)=\sum_{i}x_{i}\log x_{i}$ on the simplex.
Here $D_{h}$ is the Kullback–Leibler divergence; the algorithm becomes
the \emph{Exponentiated Gradient} method and the same $O(1/T)$ bound holds
with $D_{h}(x^{\star},x_{0})$ equal to the relative entropy between the
initial and optimal distributions.
\item \textbf{Strongly convex $f$.} If additionally $f$ is $\mu$‑strongly convex,
then a similar derivation (using the strong convexity inequality) yields
linear convergence:
\[
f(x_{k})-f(x^{\star})\le\bigl(1-\eta\mu\bigr)^{k}\,
\bigl(f(x_{0})-f(x^{\star})\bigr).
\]
\end{enumerate}

---------------------------------------------------------------------

\section*{Discussion}
The proof above follows the classic mirror‑descent analysis but
explicitly tracks every algebraic manipulation, thereby satisfying the
requirement of “no skipped steps”.  The three‑point identity (Lemma
\ref{lem:3pt}) is the cornerstone that connects the Bregman geometry to
the progress measured in the objective.  The resulting bound is
identical to that of projected Gradient Descent when the prox‑function
is Euclidean, yet the framework permits non‑Euclidean geometries that
can dramatically accelerate convergence on structured domains (e.g.,
simplex or positive orthant).  The stepsize condition $\eta\le 1/L$ is
both necessary and sufficient for the non‑positivity of the term
$\bigl(\tfrac{L}{2}-\tfrac{\sigma}{2\eta}\bigr)\|x_{k+1}-x_{k}\|^{2}$ in
\eqref{17}, ensuring the telescoping argument works.

\end{document}